\documentclass[12pt, a4paper]{article}

% --- Packages ---
\usepackage[utf8]{inputenc}
\usepackage[T1]{fontenc}
\usepackage{geometry}
\geometry{top=2.5cm, bottom=2.5cm, left=3cm, right=2.5cm}
\usepackage{amsmath, amssymb}
\usepackage{graphicx}
\usepackage{hyperref}
\usepackage{cite}
\usepackage{setspace}
\onehalfspacing

\hypersetup{
    colorlinks=true,
    linkcolor=blue,
    filecolor=magenta,      
    urlcolor=blue,
    citecolor=blue,
}

% --- Document Identity ---
\title{\textbf{Visual Recognition using Deep Learning}\\ \large Homework 2: Digit Detection Problem}
\author{
    Muhammad Rayhan Athaillah \\ 
    \small Student ID: [Your ID] \\ 
    \small National Yang Ming Chiao Tung University
}
\date{\today}

\begin{document}

\maketitle

\vspace{0.5cm}
\noindent \textbf{GitHub Repository:} \href{https://github.com/RayhanHaqi/[Your-Repo]}{https://github.com/RayhanHaqi/[Your-Repo]}

% ==========================================
% 1. INTRODUCTION
% ==========================================
\section{Introduction}
This report outlines the technical implementation for the Digit Detection Problem (Homework 2) within the Visual Recognition using Deep Learning course[cite: 4, 10]. The objective is to design a robust object detection pipeline capable of identifying and localizing multiple digits in varied RGB images[cite: 145, 148]. 

\subsection{Technical Constraints and Requirements}
The assignment dictates a specific model architecture: the system must utilize the \textbf{DEtection TRansformer (DETR)} as the core model, paired exclusively with a \textbf{ResNet-50} backbone[cite: 163]. While the backbone is fixed, modifications to other components—such as the transformer encoder and decoder—are permitted[cite: 164]. The following constraints are strictly enforced:
\begin{itemize}
    \item \textbf{Data Usage:} The use of external data from any outside source is strictly prohibited[cite: 165].
    \item \textbf{Pretrained Weights:} Pretrained weights are allowed only for the ResNet-50 backbone[cite: 164].
    \item \textbf{Submission Standards:} Final deliverables must include a ZIP file containing the source code (.py) and a PDF report written in English[cite: 21, 271, 277]. Inclusion of checkpoints or datasets in the submission is forbidden[cite: 279, 280].
    \item \textbf{Academic Integrity:} A zero-tolerance policy for plagiarism is in effect, with severe penalties including failure of the course[cite: 287, 288].
\end{itemize}

% ==========================================
% 2. LITERATURE REVIEW
% ==========================================
\section{Literature Review}

\subsection{Architectural Evolution: From DETR to RT-DETRv2}
The emergence of the DEtection TRansformer (DETR) introduced a paradigm shift by utilizing a bipartite matching loss and a transformer-based encoder-decoder to eliminate the need for hand-crafted components like Non-Maximum Suppression (NMS)[cite: 163, 164]. As proposed by Carion et al. \cite{carion2020}, DETR frames detection as a direct set prediction problem [1].

% --- Placeholder Image DETR ---
\begin{figure}[h]
    \centering
    % \includegraphics[width=0.8\textwidth]{detr_arch.png}
    \vspace{1cm} 
    \caption{Architecture of the original DEtection TRansformer (DETR) [1].}
    \label{fig:detr}
\end{figure}

Real-Time DETR (RT-DETR) later improved upon this by introducing an efficient hybrid encoder to handle multi-scale features, allowing the model to achieve speeds comparable to YOLO architectures while maintaining transformer accuracy \cite{lv2023} [2]. The latest iteration, \textbf{RT-DETRv2}, further optimizes this baseline by refining the backbone interaction and training stability, which is highly effective for dense detection tasks \cite{lv2024} [3].

% --- Placeholder Image RT-DETRv2 ---
\begin{figure}[h]
    \centering
    % \includegraphics[width=0.8\textwidth]{rtdetrv2_arch.png}
    \vspace{1cm} 
    \caption{The RT-DETRv2 architecture optimized for real-time performance [3].}
    \label{fig:rtdetrv2}
\end{figure}

\subsection{Optimization and Regularization Strategies}
Deep transformer training typically benefits from the \textbf{AdamW} optimizer \cite{loshchilov2017}, which improves generalization by decoupling weight decay from the gradient updates [4]. Within the DETR loss function, the \textbf{End-of-Sequence (\texttt{eos\_coef})} coefficient is vital for penalizing "no-object" (background) predictions. Tuning this penalty is essential to manage false positives and accommodate potential labeling noise in complex datasets.

\subsection{Resolution Adaptation}
Analysis of the provided dataset indicates a predominantly widescreen aspect ratio. Consequently, the input resolution is adapted to \textbf{$320 \times 640$} to align with the RT-DETRv2 architecture's requirements (multiples of 32) while preserving the natural geometry of the digits.

% ==========================================
% 3. METHODOLOGY
% ==========================================
\section{Methodology}

\subsection{Preprocessing and Resolution Check}
The initial preprocessing phase involved auditing the dataset resolution. To maintain spatial features without vertical distortion, images were resized to $320 \times 640$ pixels. This rectangular input ensures the model learns the correct aspect ratios of the digits.

\subsection{Data Augmentation}
To improve robustness without shifting bounding box coordinates, pixel-level augmentations were employed:
\begin{itemize}
    \item \textbf{ColorJitter:} Random adjustment of brightness, contrast, and saturation.
    \item \textbf{RandomErasing:} Strategically masking image patches to simulate oclusions.
\end{itemize}

\subsection{Regularization and Training Scheme}
Advanced regularization was implemented to ensure stable convergence:
\begin{itemize}
    \item \textbf{LR Scheduler:} The \textbf{OneCycleLR} scheduler was used to navigate the optimization landscape efficiently.
    \item \textbf{EMA:} \textbf{Exponential Moving Average} of model weights was utilized to enhance the stability of final predictions.
\end{itemize}

\subsection{Hyper-parameter Experiments}
Systematic experiments were conducted using a Dual NVIDIA RTX 4090 setup, focusing on the following variables:
\begin{itemize}
    \item \textbf{Object Queries ($q$):} Evaluation of $q=1000$ versus $q=1500$ to accommodate digit density.
    \item \textbf{Weight Decay (WD):} Comparative testing between $1 \times 10^{-3}$ and $5 \times 10^{-4}$ for optimal regularization.
    \item \textbf{Learning Rate (LR):} Tuning the base learning rate at $1 \times 10^{-4}$ to balance speed and accuracy.
\end{itemize}

% ==========================================
% REFERENCES
% ==========================================
\newpage
\begin{thebibliography}{99}
\bibitem{carion2020} M. Carion et al., ``End-to-End Object Detection with Transformers,'' in \textit{Proc. ECCV}, 2020.
\bibitem{lv2023} W. Lv et al., ``DETRs Beat YOLOs on Real-time Object Detection,'' in \textit{Proc. CVPR}, 2023.
\bibitem{lv2024} W. Lv et al., ``RT-DETRv2: Improved Baseline with Faster Training and Better Generalization,'' \textit{arXiv:2407.17140}, 2024.
\bibitem{loshchilov2017} I. Loshchilov and F. Hutter, ``Decoupled Weight Decay Regularization,'' in \textit{Proc. ICLR}, 2019.
\end{thebibliography}

\end{document}